\section{Hardware Counters}
\label{section:core:hardware-counters}

Core performance is a question of how well we are exploiting the single-core
hardware features on a chip. Thankfully hardware vendors typically include
extra registers on the chip which record the instance of certain hardware
related events. These added registers are known as hardware performance
counters and can be used to diagnose performance issues of a code running on a
specific architecture. For example, cache misses can be registered on cache
related hardware counters and so can investigate whether a code has inefficient
use of the cache heirarchy and is hence having to increase data traffic from
memory. Other hardware counters can activate through branch mispredictions and
so we can target control flow within a code that is hampering performance.

One challenge of working with hardware counters is that they are not
standardised across all chip generations and vendors. Hence, an analyst has to
be relatively flexible when using hardware counters. 

\subsection{Accessing hardware counters}

There are multiple tools built for accessing hardware counters and other tools
which can use hardware counters to integrate into other metrics for analyses.
We do not delve in to the details here, but there are many performance tools
such as Score-P, Scalasca and Cube which offer lots of functionality with code
instrumentation, profiling and trace analysis that can also access hardware
counters often through the PAPI (Performance Application Programming Interface)
library. More information of accessing hardware counters in Score-P is given in
Section~\ref{section:localisation:scorep}.

On many Linux systems there is the \texttt{perf} which is quite a general tool
for performance though can be relatively complex to use. However, as a quick
aside, \texttt{perf} can be used to record metrics for Score-P. A more
HPC-focused command line tool is the \texttt{likwid-perfctr} tool from the
LIKWID suite. Rather than giving just simply read outs of hardware counters,
which can be very difficult to interpret for beginners, \texttt{likwid-perfctr}
aggregates the hardware counter results into easier to interpret metrics.

With \texttt{likwid-perfctr} we have already how to gather data from
performance groups such as \texttt{FLOPS\_DP} for double-precision floating
point operations. 

\begin{wip}
Do we include how to gather hardware counter data with vendor tools like AMD
	uProf and Intel's VTune?
\end{wip}

\subsection{Gathering the data}
\begin{wip}
Very early work here, likely to include some pointers on a simple workflow for
	gathering multiple runs with LIKWID
\end{wip}

There can be significant variability in some hardware counters and so we advise
doing batches of runs to aggregate statistics.
